\begin{abstract}
We study CounterTrace Necessity Training (CNT) as an object-first approach to reasoning supervision. The target object is trajectory-conditional necessity: whether a local step remains necessary for successful continuation once nearby edits and continuators are controlled. Rather than presenting a training-win story, we report a staged empirical status update. First, the object is recovered on synthetic data, where CNT tracks programmed ground truth with correlation $0.9902$ and does not collapse onto observational or entropy-style shortcuts. Second, the same object survives a larger-pool audited GSM8K pipeline: from a 112-trace verified collection baseline, 230 of 298 candidate step pairs survive the final joint audit. Third, current offline Stage~B families still fail to convert that object signal into a stable rollout utility gain over matched SFT-only control. The supported claim is therefore narrower but useful: CNT yields a reproducible benchmark-and-audit protocol for necessity, while offline utility conversion remains unresolved.
\end{abstract}
